% Pista geo-23j-q6 · Geometria
% Procedència: PAU Catalunya 2023 · Convocatòria de juny · Sèrie 1 · Qüestió 6
\documentclass[11pt,a4paper]{article}
\usepackage{pista}
\pistainfo{Catalunya 2023}{Convocatòria de juny}{Sèrie 1}

\begin{document}

\section*{\textcolor{blauPista}{Qüestió 6}}

\noindent Siguin els plans $\pi_{1}$ i $\pi_{2}$, determinats respectivament per les equacions $\pi_{1}: x + y = 3$ i $\pi_{2}: x - z = -2$.

\subsection*{\textit{a)} \normalfont Trobeu l'equació general ($Ax + By + Cz + D = 0$) del pla $\pi_{3}$, que és perpendicular a $\pi_{1}$ i $\pi_{2}$, i que passa pel punt $P = (4, 1, 2)$. \hfill\textnormal{[0,75~punts]}}

\pista{Recorda que dos plans són perpendiculars quan els seus vectors normals també ho són. Per tant, perquè $\pi_{3}$ sigui perpendicular alhora a $\pi_{1}$ i a $\pi_{2}$, el seu vector normal $\vec{n}_{3}$ ha de ser perpendicular alhora a $\vec{n}_{1}$ i a $\vec{n}_{2}$. Quina operació entre dos vectors et dona directament un vector perpendicular als dos? Un cop hagis obtingut $\vec{n}_{3}=(A,B,C)$, gairebé ho tens tot, perquè pots escriure l'equació de $\pi_{3}$ de la forma $Ax + By + Cz + D = 0$ i només et faltarà trobar el valor de $D$. Per trobar $D$, n'hi ha prou de substituir les coordenades de $P$ a l'equació de $\pi_{3}$.}

\subsection*{\textit{b)} \normalfont Sigui $r$ la recta d'intersecció de $\pi_{1}$ i $\pi_{2}$. Calculeu l'equació vectorial de la recta $r$. \hfill\textnormal{[0,75~punts]}}

\pista{La recta $r$ és la intersecció dels dos plans, així que està continguda alhora en tots dos. Per tant, la seva direcció és perpendicular alhora a $\vec{n}_{1}$ i a $\vec{n}_{2}$. Fixa't que aquesta és, exactament, la mateixa direcció que has trobat a l'apartat a) per al vector normal de $\pi_{3}$: ja la tens! Per poder trobar un punt que pertany a la recta $r$, n'hi ha prou amb trobar una solució del sistema format per les dues equacions dels plans (per exemple, decidint el valor $z = 0$ i resolent les dues equacions resultants per a $x$ i $y$). Amb un punt i un vector director, obtenir l'equació vectorial de la recta és immediat.}

\subsection*{\textit{c)} \normalfont Calculeu el punt $Q$ de la recta $r$ que és més a prop del punt $P$. \hfill\textnormal{[1~punt]}}

\pista{El punt de $r$ més proper a $P$ és el peu de la perpendicular des de $P$ a la recta. Fixa't que el pla $\pi_{3}$ de l'apartat a) passa per $P$ i és perpendicular a $r$: per tant, $Q$ és la intersecció de $r$ amb $\pi_{3}$. Parametritza els punts de $r$ amb el paràmetre $t$ fent servir l'equació vectorial de l'apartat b), escriu el vector $\vec{PQ}$ en funció de $t$, i imposa que $Q$ sigui de $\pi_{3}$: com que el vector director $\vec{v}_{r}$ de $r$ és normal a $\pi_{3}$, això vol dir $\vec{PQ} \cdot \vec{v}_{r} = 0$. Obtindràs una equació en $t$; un cop sàpigues quant val $t$, retorna a la parametrització de $r$ per obtenir les coordenades concretes del punt $Q$.}

\end{document}
